% TOP_PROBLEM: {"id": 3418, "title": "Is there an algorithm to determine if a triangulated 4-manifold is combinatorially equivalent to the 4-sphere?", "category_id": 7, "category": {"id": 7, "name": "topology", "display_name": "Topology", "description": "Properties preserved under continuous deformations.", "slug": "topology", "order_index": 7, "created_at": "2026-07-31T15:26:25.671Z"}, "status": "open", "problem_number": "OPG-37161", "set_id": 12}
% FIRST_POSTED: 2026-10-01
% ENTRY_KIND: statement_only
% UnsolvedMath ID 3418; source catalogue code OPG-37161
% !TeX program = lualatex
% Definition and sources only; this file is not an LLM solution attempt.
\documentclass[11pt,a4paper]{article}
\usepackage[margin=25mm]{geometry}
\usepackage{fontspec}
\setmainfont{lmroman10-regular.otf}[BoldFont=lmroman10-bold.otf,
 ItalicFont=lmroman10-italic.otf,BoldItalicFont=lmroman10-bolditalic.otf]
\usepackage{amsmath,amssymb,mathtools}
\usepackage{unicode-math}
\setmathfont{latinmodern-math.otf}
\AtBeginDocument{\let\setminus\smallsetminus}
\usepackage{xurl,hyperref}
\hypersetup{colorlinks=true,urlcolor=blue}
\setcounter{secnumdepth}{0}
\setlength{\parindent}{0pt}
\setlength{\parskip}{5pt}
\setlength{\emergencystretch}{3em}
\begin{document}
\section*{Is there an algorithm to determine if a triangulated 4-manifold is combinatorially equivalent to the 4-sphere?}\label{3418}
{\small UnsolvedMath ID 3418 · OPG-37161 · Topology · OpenGarden}
\subsection{Definitions and mathematical statement}
The input is a finite simplicial complex $K$, presented by its finite
vertex and simplex lists, promised to triangulate a closed connected
piecewise-linear four-manifold. In particular its vertex links are
piecewise-linearly homeomorphic to $S^3$. Let $\partial\Delta^5$ be
the boundary complex of the five-simplex, a standard triangulation
of the four-sphere.

Is there an algorithm that halts on every promised input and decides
whether
\[
                         |K|\cong_{\mathrm{PL}}|\partial\Delta^5|?
\]
Here PL equivalence means that after suitable finite subdivisions
the triangulated spaces admit a piecewise-linear homeomorphism.
It does not mean that the original simplex lists are isomorphic.
Equivalently one can use existence of a finite sequence of bistellar
(Pachner) moves and simplicial isomorphisms joining the triangulations.

The requested procedure has no time-complexity bound, but it must
terminate correctly in both cases. Enumerating successful move
sequences only recognizes yes-instances and supplies no termination
test on no-instances. The input promise avoids conflating recognition
of a four-manifold with recognition of the standard PL sphere.
The category remains PL throughout: topological homeomorphism or
homotopy equivalence alone is not the specified equivalence.
The source's discussion of smoothings is unnecessary for stating
this exact finite-input decision problem.
\subsection{Short English statement}
Is there an always-terminating algorithm recognizing the standard piecewise-linear four-sphere among finite triangulated closed four-manifolds?
\subsection{Sources}
\begin{itemize}
\item[S1] ULAM.AI and the UnsolvedMath Contributors, supplied problems.json, record 3418 (OPG-37161), OpenGarden collection. Catalogue identity and source transcription; CC BY 4.0.
\url{https://huggingface.co/datasets/ulamai/UnsolvedMath}.
\item[S2] Open Problem Garden, Is there an algorithm to determine if a triangulated 4-manifold is combinatorially equivalent to the 4-sphere?. Original problem and discussion; the mathematical formulation below is expanded from the supplied source record.
\url{http://www.openproblemgarden.org/op/is_there_an_algorithm_to_determine_if_a_triangulated_4_manifold_is_combinatorially_equivalent_to_the_4_sphere}.
\end{itemize}
\end{document}
